\documentclass[sigconf]{acmart}
% Original report text. Template structure adapted from the MIT reference
% repository at 81a739bb7786a44b989f5d76afbb6f9d22b5d1dd.
% No source-report prose, author details, figures, or results are copied.
% Non-ACM competition report metadata; ACM formatting does not imply ACM publication.
% Structural setup adapted from the MIT reference repository, commit
% 81a739bb7786a44b989f5d76afbb6f9d22b5d1dd, reports/preamble.tex.
\setcopyright{none}
\settopmatter{printacmref=false}
\acmConference[Kaggle Competition]{Filament Segmentation Challenge}{2026}{\href{https://www.kaggle.com/competitions/filament-segmentation-2026}{kaggle.com/competitions/filament-segmentation-2026}}
\acmBooktitle{Competition Report}
\acmDOI{}
\acmISBN{}
\usepackage{mathtools}
\usepackage{tabularx}
\newcommand{\pending}{\textnormal{pending}}
\raggedbottom

\begin{document}
\title[Temporal Validation for Solar Filament Segmentation]{Solar Filament Instance Segmentation with Temporal Validation and Resolution-Aware Diagnostics}
\subtitle{A Solution to the Solar Filament Segmentation Challenge 2026}
\author{Xin Meng}
\affiliation{\institution{China Medical University}\country{China}}
\email{mengxinlab@gmail.com}
\renewcommand{\shortauthors}{Xin Meng}
\begin{abstract}
We study image-only instance segmentation of solar filaments in 2048$\times$2048 grayscale GONG H-alpha observations. A reproducible local-contrast baseline achieved development Panoptic Quality (PQ) 0.177047 and, after selection was frozen, PQ 0.105604 on 49 later observations. A compact whole-frame U-Net and a fixed neural--contrast mixture did not improve the development score. Diagnostics identify low neural responses, sparse positive pixels, and loss of filament scale under whole-frame downsampling, without evidence of input--target misalignment. A prespecified native-resolution V4 comparison restored nonzero matches in all development periods but did not improve the reference. An empty-ring defect invalidated the normalization arm; it is retained as a failed experiment, not evidence about normalization efficacy. The first formal Kaggle submission was accepted with Public Score 0.15; this score verifies submission behavior and is excluded from parameter selection.
\end{abstract}
\maketitle
\noindent\textbf{Evidence cutoff, 30 October 2026:} v3 is frozen; V4 raw comparisons are valid, normalization comparisons are invalid, and further GPU spending fails the prespecified gate.

\section{Problem and Data}
Solar filaments are dark, elongated structures with weak contrast, fine barbs, fragmented visible material, and nearby non-filament dark regions. The competition requires one mask per filament, making instance identity distinct from semantic foreground detection~\cite{competition}. We seek reproducible image-only inference and a validation protocol that exposes temporal failures.

The supplied archive contains 707 training and 180 test observations, all 8-bit grayscale JPEGs at 2048$\times$2048. Its COCO-style training file has 1,154 annotator--image records and 8,199 instance annotations. Independent annotators can label the same physical observation. We use only the competition training segmentation polygons. Spines, boxes, chirality labels, annotator identity, and date/station metadata are not predictive features; dates are used for partitioning.

External MAGFiLO annotations are excluded because the host restricts annotation supervision to the supplied training labels and public discussion reports overlap with test imagery. We do not resolve test IDs against external labels, use other submissions, manually label test data, or exploit metric edge cases. A supplied-data exact-byte audit found no duplicate files. Some low-resolution similarity flags were attributable to common solar-disk shape when flagged training pairs were examined. This does not prove the absence of every possible dependency.

\section{Temporal Protocol and Official PQ}
The source-image blocks are 2011--2013, 2014--2016, 2017--2019, and 2020--2022, containing 279, 268, 111, and 49 images. All annotators of an observation remain together. The first three blocks are development folds; training observations within 27 days of validation dates are purged. Development also excludes observations within 27 days of the fourth block. Exact memberships are archived.

The v3 method was selected on development data before its single frozen assessment on the fourth block. V4 excludes that block's labels from development. If they are used later, that subset must be explicitly downgraded to development, and later scores cannot be described as untouched holdout evidence. Public reference parameters already reflect their author's competition-data tuning, limiting claims of pristine validation.

The host notebook evaluates each annotator--image separately against the same image prediction set~\cite{selfeval}. Every pair with strict $\operatorname{IoU}>0.5$ is a TP; the code does not use one-to-one assignment. FP and FN are zero-hit prediction columns and GT rows. Counts and hit IoUs are accumulated before computing
\begin{equation}
 \operatorname{PQ} =
 \frac{\sum_{(g,p):\,\operatorname{IoU}(g,p)>0.5}\operatorname{IoU}(g,p)}
 {\operatorname{TP}+0.5\operatorname{FP}+0.5\operatorname{FN}}.
\end{equation}
Mean image PQ and annotator-unioned GT would change this protocol. A same-family provisional audit found semantic agreement with our RLE evaluator, with a minor float32/float64 difference: a $2/3$-IoU witness differed by $1.99\times10^{-8}$, without changing hit counts. Frozen v3 scores are not rewritten. Any numeric-compatibility evaluator has its own version.

\clearpage
\section{Frozen v3 Methods and Results}
\subsection{Contrast and Neural Models}
The selected contrast baseline adapts Sijan Regmi's MIT code~\cite{regmi}. It finds a bright solar disk using Otsu thresholding, the largest connected region, contour filling, and erosion. Relative darkness is computed from a smoothed image and Gaussian local background. Selected parameters are background scale 64 native pixels, threshold 0.04, no closing, component area 800--100,000 native pixels, and minimum elongation 2. Disjoint components become instances. These are a public reference, not a novel method.

Our compact U-Net~\cite{unet} is trained from scratch with grayscale input, GroupNorm, SiLU, and widths 16/32/64/128. Each entire observation is resized to 512$^2$. Independent annotator foreground unions are averaged into a soft target and area-resized. This union is used only for training: evaluation retains every original 2048$^2$ instance and annotator separately.

Training uses BCE plus soft Dice, AdamW, fixed seeds, and matched image/target dihedral transformations. Eight epochs of 48 updates at batch size four yield 384 optimizer updates and 1,536 sampled whole-frame draws per fold, not eight exhaustive dataset passes. There is no foreground-directed sampling or neural disk-intensity normalization. Inference averages identity, horizontal, vertical, and combined flips, restores probabilities to 2048$^2$, masks off-disk regions, and extracts thresholded components. Four predefined threshold/closing variants and a fixed 50:50 neural--contrast blend were compared.

\begin{table}[t]
\caption{Selected v3 development candidates; all eleven configurations and period results are archived.}
\label{tab:v3}
\centering
\begin{tabular}{lrrrr}
\toprule
Candidate & PQ & TP & FP & FN\\
\midrule
Contrast reference & .177047 & 2245 & 6084 & 5679\\
50:50 hybrid & .173176 & 2231 & 6504 & 5693\\
Best pure U-Net & .018415 & 129 & 527 & 7795\\
\bottomrule
\end{tabular}
\end{table}

\subsection{Results and Formal Submission}
The reference's period PQ values are 0.174206, 0.177008 and 0.198192. The hybrid added 420 FP and lost 14 TP; its small increase in hit-IoU sum did not offset the larger denominator. The pure neural model's period PQ was 0, 0 and 0.137395. These results selected the contrast reference.

Its single later-block assessment yielded PQ \textbf{0.105604}, TP/FP/FN \textbf{49/261/226}. The absolute drop from development is 0.071443, or 40.4\%. A previously calculated calendar-month block-bootstrap interval is [0.075674, 0.127158], conditional on the selected model and observed holdout months. It is not hidden-test score uncertainty. Temporal generalization remains weak.

The frozen CSV contains 1,312 nonempty instances on 175 test images; five images naturally have zero detections. Their IDs remain in the coverage manifest. No dummy masks were inserted. The file was formally uploaded as Kaggle submission \textbf{57795106}, accepted with status \textbf{COMPLETE} and Public Score \textbf{0.15}. The canonical receipt is \texttt{docs/official\_submission\_v3.md}. This score is a production submission check and is not used to choose models, thresholds, or fusion parameters.

Each row has a unique \texttt{filament\_id} formed from the original image stem plus an instance suffix, and \texttt{segmentation\_rle} containing compressed ASCII COCO counts. Masks are binary Fortran-order uint8 arrays at 2048$^2$. The frozen CSV SHA-256 begins \texttt{c8e6ac199a3fbe42}; the full value is recorded with the submission. Public Score and local PQ refer to different sample sets and are reported separately.

The completed run used Tesla T4 (14.6 GiB reported VRAM), Python 3.13.15, PyTorch 2.11.0+cu128, and CUDA 12.8. Neural training phases took 223.9, 218.0, and 222.0 seconds, excluding preprocessing, classical search, validation inference, and packaging. No paid resource was provisioned.

\clearpage
\section{Neural Failure Diagnosis}
The diagnostic analysis uses development observations and cached predictions; it does not inspect the frozen holdout's masks to develop V4. Existing v3 evidence remains immutable.

\begin{table}[t]
\caption{Development diagnostics. Maxima bound all validation maps in the indicated fold; foreground fractions and widths come from the prespecified diagnostic sample.}
\label{tab:diag}
\centering
\begin{tabular}{lrrr}
\toprule
Quantity & Fold 0 & Fold 1 & Fold 2\\
\midrule
Soft FG fraction & .003108 & .003369 & .000872\\
Native median width (px) & 9.297 & 8.549 & 7.928\\
Width after $4\times$ resize & 2.32 & 2.14 & 1.98\\
PQ at threshold .60 & 0 & 0 & .137395\\
\bottomrule
\end{tabular}
\end{table}

All 268 fold-1 validation maps have maximum probability at most 0.154663, below the lowest tested threshold of 0.30; zero instances and zero TP follow from that cutoff. Fold-0 maxima are at most 0.503418, so a 0.60 threshold produces no instances. Lowering fold 0 to 0.30 still gives only 6 TP, 201 FP, and 3,725 FN. Low neural response is established; a permissive threshold is not established as a solution.

On 18 prespecified development images, reconstructed cached targets agree within 0.001960814, the expected half-step uint8 quantization error. Image and target augmentations correspond. These checks provide no evidence of an alignment defect. The canonical diagnosis is \texttt{diagnostics/v4/p2\_unet\_diagnosis.md} and its JSON companion.

Native median widths of roughly eight to nine pixels become only two pixels after fourfold resizing (Table~\ref{tab:diag}). Seventeen of 155 sampled instances have hard 512$^2$ areas below 25 pixels, the resized equivalent of the 400-native-pixel component filter. Whole-frame resizing and component rejection can therefore affect thin structures. Foreground is sparse, and the neural sampling policy does not target it.

The combined evidence motivates native-resolution supervision, explicit sampling, and intensity normalization. It does not separately identify each factor's causal effect. Declining training loss cannot establish overfitting because the histories contain no validation learning curves. Blindly increasing epochs is not the first proposed intervention.

\section{Method Reference and Controlled V4 Design}
We inspect the MIT repository \texttt{ameyypawar/filament-seg}, commit \texttt{81a739bb7786}~\cite{pawar}; the full hash is in the reference. It uses native 512$^2$ crops, foreground-directed sampling, radial median-profile limb-darkening correction, corrected intensity plus radius channels, and an ImageNet-initialized ResNet34 U-Net~\cite{resnet}. Its detector fusion assigns unclaimed semantic foreground to selected, dilated YOLO11-seg regions.

Its 563/144 split randomly groups dates without our embargo; postprocessing cross-fitting does not make the complete research process a blind temporal evaluation. Its greedy one-to-one scorer and treatment of prediction-only images differ from the host notebook. We neither inherit that scorer as official PQ nor compare its validation score directly to our holdout.

V4 freezes two input arms, \texttt{raw\_radius} and \texttt{flat\_radius}, under the same network, initialization, and update budget. Both use native 512$^2$ crops, 50\% foreground-directed sampling, 384 updates per fold, and 25 native-resolution inference tiles per image. The only input-arm intervention is radial intensity correction; radius is available in both. Evaluation compares connected components (\texttt{cc}) and a fixed 6-pixel component-grouping option (\texttt{group6}). Threshold 0.45 and minimum area 400 native pixels are fixed before results. Exact sampling, normalization, reconstruction, and grouping definitions are retained in the frozen configuration and code.

Overlapping tile predictions are consolidated before instance extraction. No duplicated crop masks, nested copies, GT-informed instance counts, or leaderboard parameter scans are permitted. All original development periods and the 27-day embargo are retained. The fourth block is excluded from V4 development.

\clearpage
\section{V4 Evidence and GPU Decision}
\begin{table}[t]
\caption{V4 development results. Flat rows ($^\ast$) are \textbf{INVALID}: recorded scores are retained solely as failure evidence.}
\label{tab:v4}
\centering
\small
\setlength{\tabcolsep}{3pt}
\begin{tabular}{llrrrr}
\toprule
Input & Instances & PQ & TP & FP & FN\\
\midrule
raw & cc & .119969 & 1519 & 6369 & 6405\\
raw & group6 & .127735 & 1614 & 6216 & 6310\\
flat$^\ast$ & cc & .127519 & 2042 & 10518 & 5882\\
flat$^\ast$ & group6 & .139628 & 2261 & 10382 & 5663\\
\bottomrule
\end{tabular}
\end{table}
\begin{table}[t]
\caption{Period PQ under the unchanged date-grouped development protocol. Starred values remain invalid recorded outputs.}
\label{tab:period}
\centering
\small
\setlength{\tabcolsep}{3pt}
\begin{tabular}{llrrr}
\toprule
Input & Instances & 2011--13 & 2014--16 & 2017--19\\
\midrule
raw & cc & .144486 & .104235 & .080279\\
raw & group6 & .155649 & .108041 & .086687\\
flat$^\ast$ & cc & .129835 & .132953 & .096548\\
flat$^\ast$ & group6 & .141781 & .144623 & .109522\\
\bottomrule
\end{tabular}
\end{table}
The native-resolution raw model recovers valid matches in all periods, unlike the two zero-match v3 neural folds. It does not improve the classical reference. Grouping raises raw PQ from 0.119969 to 0.127735, with 95 additional TP and 153 fewer FP, but the best valid V4 result is 27.9\% below development PQ 0.177047. Each of its period scores also remains below the corresponding v3 reference score.

\paragraph{Invalid normalization experiment.}
The original radial-profile implementation resized a disk sample to 256$^2$ and divided it into 128 rings. Observation \texttt{20110918112834Ch} has no pixels in ring 127; its median becomes NaN. Cache conversion to uint8 writes zero at non-finite entries while floating-point inference retains NaNs, producing an invalid train/inference preprocessing contract. Thus both flat rows in Tables~\ref{tab:v4}--\ref{tab:period} are invalid, regardless of recorded PQ. They cannot establish whether limb-darkening correction helps or harms. Original outputs are preserved. The correction estimates all 128 ring medians from native 2048$^2$ disk pixels; a single real-case CPU regression found at least 137 pixels per ring and finite, positive profile values, without empty-bin filling. It passed in 0.0776 seconds. No GPU model was retrained after this fix; no repaired-arm score exists. Evidence is retained in \path{diagnostics/v4/p4_radial_failure.md} and its JSON companion.

\paragraph{Compute and decision.}
The run took 3,427.583 seconds (57.13 minutes) on T4, within the two-hour ceiling. Cache preparation took 339.101 seconds. Raw training took 650.989 seconds and inference 791.166 seconds; flat training took 638.814 seconds and inference 992.412 seconds. Peak allocated CUDA memory was 2.084 GiB per arm.

The prespecified continuation gate requires development PQ at least 0.197047225, improvement in at least two periods, and minimum period PQ greater than 0.05. The best valid candidate fails the pooled-score and period-improvement requirements. The decision is therefore \textbf{STOP additional GPU spending and retain the accepted v3 model}. Invalid flat scores are excluded from candidate selection and gate adjudication. A repaired normalization comparison would require a new frozen plan rather than a silent retry, and the Public Leaderboard is not a tuning signal. Canonical V4 outputs remain separate under \path{results/kaggle_v4/project/runs/v4_native}.

\paragraph{Illustrative raw-model failures.}
Observation \texttt{20120703070814Th} has 0 TP, 24 FP and 43 FN: weak, long filaments are missed or fragmented. Observation \texttt{20171029150150Bh} has 0 TP, 51 FP and 15 FN, with many false regions along the solar limb. These counts pool all annotators; the archived overlays show one declared annotator. Cases were selected from existing error counts, with no further parameter tuning.

\section{Reproducibility, Licensing and Limits}
Versioned source, configurations, image memberships, histories, weights, predictions, PQ counts, environments, and receipts are preserved. v3 artifacts remain under \path{results/kaggle_v3/project/}; V4 is separate, including invalid outputs and adjudication. Fixed seeds do not imply bit-identical results across devices or library versions.

The \href{https://github.com/mengxinlab/solar-filament-2026-reproducible}{public Git repository} and \href{https://www.kaggle.com/code/mengxinlab/solar-filament-2026-reproducible-release}{Kaggle notebook} supply pinned dependencies, complete pipeline code, selected classical parameters and inference without retraining. Raw competition data remain on Kaggle under their CC BY-NC terms.

MIT notices and attribution accompany the reference code. They do not relicense dependencies or weights; Ultralytics is separately AGPL-3.0. We do not download the studied author's competition-trained weights, which could have used our holdout images. V4 uses no external MAGFiLO labels. Ordinary Kaggle GPU is used; ARC-reserved RTX6000 is excluded.

Absolute PQ remains modest. The raw bundle changes resolution and sampling jointly, without isolating their individual effects. Normalization efficacy remains untested. We do not claim fine-barb accuracy, future-cycle robustness, or independent validation; routine review is same-family and provisional.

The separate release checklist records final form fields and artifact hashes. No Google Form is submitted by this report.

\bibliographystyle{ACM-Reference-Format}
\bibliography{main}
\end{document}
